% Generated by support/export-artifacts.py; do not edit.
% Slug: ch40-backward
% Source: chapters/ch40.tex:284
% TikZ source: figures/logic-backward-assertions.tex:1
% Original label: \label{fig:ch40-backward}
% Original caption (LaTeX source):
%   Execution moves down the command sequence while substitution moves the obligation up it.
% Caption: Execution moves down the command sequence while substitution moves the obligation up it.
% Setup: support/compiler-pseudocode.sty (TikZ libraries and shared styles).
\begin{tikzpicture}[node distance=14mm]
\node[basic block] (first) {$a:=a+2$};
\node[basic block, below=of first] (second) {$b:=3a$};
\node[basic block, below=of second] (end) {finish};
\draw[cfg edge] (first) -- (second);
\draw[cfg edge] (second) -- (end);
\node[right=20mm of first, align=center] (p) {$3(a+2)\geq12$};
\node[right=20mm of second, align=center] (r) {$3a\geq12$};
\node[right=20mm of end, align=center] (q) {$b\geq12$};
\draw[cfg edge] (q) -- (r);
\draw[cfg edge] (r) -- (p);
\node[above=5mm of first, font=\small] {execution};
\node[above=5mm of p, font=\small] {obligation};
\end{tikzpicture}

% Original figure LaTeX:
% \begin{figure}
% \centering
% \begin{tikzpicture}[node distance=14mm]
\node[basic block] (first) {$a:=a+2$};
\node[basic block, below=of first] (second) {$b:=3a$};
\node[basic block, below=of second] (end) {finish};
\draw[cfg edge] (first) -- (second);
\draw[cfg edge] (second) -- (end);
\node[right=20mm of first, align=center] (p) {$3(a+2)\geq12$};
\node[right=20mm of second, align=center] (r) {$3a\geq12$};
\node[right=20mm of end, align=center] (q) {$b\geq12$};
\draw[cfg edge] (q) -- (r);
\draw[cfg edge] (r) -- (p);
\node[above=5mm of first, font=\small] {execution};
\node[above=5mm of p, font=\small] {obligation};
\end{tikzpicture}

% \caption{Execution moves down the command sequence while substitution moves the obligation up it.}
% \label{fig:ch40-backward}
% \end{figure}
